\section*{第 \ref{ch:b1:seq} 章——数列}

\begin{solution}{exo:b1:seq:1}
$\Bigl|\dfrac{2n+1}{n+3} - 2\Bigr| = \dfrac{5}{n+3}$。给定
$\varepsilon > 0$，由\omterm{thm:b1:reals:archimedes}{阿基米德性质}取 $N > \frac 5\varepsilon - 3$：
当 $n \geq N$ 时，$\frac{5}{n+3} \leq \varepsilon$。于是极限为
$2$。

$((-1)^n)$：其\omterm{def:b1:seq:subsequence}{子列} $(u_{2n}) = (1)$ 与 $(u_{2n+1}) =
(-1)$ 收敛于不同的极限，因此该数列发散
（见\cref{prop:b1:seq:subsequences}）。（直接验证：任取候选极限 $\ell$
对 $\varepsilon = \frac12$ 不满足定义，因相邻两项的距离为
$2$。）
\end{solution}

\begin{solution}{exo:b1:seq:2}
两边同除以 $n^2$：$\dfrac{1 - 3/n + 1/n^2}{2 + 5/n^2} \to
\dfrac12$。

$\sqrt{n+1} - \sqrt n = \dfrac{1}{\sqrt{n+1} + \sqrt n} \to 0$
（利用共轭式）。

$\dfrac{2^n + n^3}{3^n - n^2} = \dfrac{(2/3)^n + n^3/3^n}{1 -
n^2/3^n} \to \dfrac{0 + 0}{1 - 0} = 0$，利用 $q^n \to 0$（$\abs q
< 1$ 时），以及\omterm{def:b1:poly:def}{多项式}与等比数列的比较
（见\cref{prop:b1:functions:powerrules}）。

$5^{1/n} = \eu^{\frac{\ln 5}{n}} \to \eu^0 = 1$。
\end{solution}

\begin{solution}{exo:b1:seq:3}
伯努利不等式：$(1+h)^n \geq 1 + nh$，其中 $h \geq -1$，用归纳法证明——
$(1+h)^{n+1} = (1+h)^n(1+h) \geq (1+nh)(1+h) = 1 + (n+1)h + nh^2
\geq 1 + (n+1)h$。

当 $0 < \abs q < 1$ 时，记 $\frac{1}{\abs q} = 1 + h$，$h > 0$，则
$\abs{q}^n = \frac{1}{(1+h)^n} \leq \frac{1}{1 + nh} \to 0$，由夹逼准则
可得 $q^n \to 0$（$q = 0$ 时平凡）。对 $q = 1$：
为常数列，极限为 $1$。对 $q = -1$，数列发散
（见\cref{exo:b1:seq:1}）。对 $\abs q > 1$，$\abs q^n = (1 + h)^n \geq
1 + nh \to +\infty$，故 $(q^n)$ 无界，从而发散（
趋于 $+\infty$，若 $q > 1$；若 $q <
-1$，则符号交替，无极限）。
\end{solution}

\begin{solution}{exo:b1:seq:4}
不动点：由 $\ell = \frac{\ell + 3}{2}$ 得 $\ell = 3$。于是
\[
v_{n+1} = u_{n+1} - 3 = \frac{u_n + 3}{2} - 3 = \frac{u_n - 3}{2} =
\frac{v_n}{2}:
\]
$(v_n)$ 是公比为 $\frac12$ 的等比数列，$v_0 = -3$。于是 $u_n = 3 -
\frac{3}{2^n} \to 3$。
\end{solution}

\begin{solution}{exo:b1:seq:5}
$H_{2n} - H_n = \sum_{k=n+1}^{2n} \frac 1k \geq n \cdot
\frac{1}{2n} = \frac12$（这 $n$ 项中每项都 $\geq
\frac{1}{2n}$）。假设 $(H_n)$ 为\omterm{def:b1:seq:cauchy}{柯西序列}，取 $\varepsilon =
\frac13$ 时必有 $\abs{H_{2n} - H_n} \leq \frac13$，对充分大的
$n$ 成立，矛盾。单调递增而不收敛的数列发散到
$+\infty$（见\cref{thm:b1:seq:monotone}）：$H_n \to +\infty$。
\end{solution}

\begin{solution}{exo:b1:seq:6}
设 $a = \lim u_{2n}$，$b = \lim u_{2n+1}$，$c = \lim u_{3n}$。
数列 $(u_{6n})$ 既是 $(u_{2n})$ 的\omterm{def:b1:seq:subsequence}{子列}，也是
$(u_{3n})$ 的\omterm{def:b1:seq:subsequence}{子列}，其极限同时为 $a$ 与 $c$，故 $a = c$。数列
$(u_{6n+3})$ 是 $(u_{2n+1})$ 的\omterm{def:b1:seq:subsequence}{子列}（下标为奇数），也是
$(u_{3n})$ 的\omterm{def:b1:seq:subsequence}{子列}（下标为 $6n + 3 = 3(2n+1)$），故 $b = c$。于是 $a = b$，
由\cref{prop:b1:seq:subsequences}（偶数项与奇数项的极限
相等）可知 $(u_n)$ 收敛。
\end{solution}

\begin{solution}{exo:b1:seq:7}
\emph{稳定性与界：}$I = \intcc{0}{2}$ 是稳定的：对 $x \in
I$，$\sqrt{2 + x} \in \intcc{\sqrt 2}{2} \subseteq I$；且 $u_0 = 0
\in I$。

\emph{单调性：}$f(x) = \sqrt{2+x}$ 单调递增，且 $u_1 =
\sqrt 2 > u_0$，由归纳法得 $(u_n)$ 单调递增，
且有\omterm{def:b1:reals:bounds}{上界} $2$，因此
由\cref{thm:b1:seq:monotone}可知它收敛。

\emph{极限：}由 $\ell = \sqrt{2 + \ell}$，$\ell \geq 0$，得
$\ell^2 - \ell - 2 = 0$，故 $\ell = 2$。

\emph{误差估计：}乘以共轭式，
\[
2 - u_{n+1} = 2 - \sqrt{2 + u_n}
= \frac{4 - (2 + u_n)}{2 + \sqrt{2+u_n}}
= \frac{2 - u_n}{2 + \sqrt{2 + u_n}}
\leq \frac{2 - u_n}{3},
\]
由于 $\sqrt{2 + u_n} \geq \sqrt 2 > 1$。由归纳法，从 $2 - u_0
= 2$ 出发，可得 $\;0 \leq 2 - u_n \leq \frac{2}{3^n}$。
\end{solution}

\begin{solution}{exo:b1:seq:8}
假设 $(u_n)$ 不收敛于 $\ell$：存在 $\varepsilon_0
> 0$，使得无穷多个下标满足 $\abs{u_n - \ell} >
\varepsilon_0$，这些下标构成\omterm{def:b1:seq:subsequence}{子列} $(u_{\varphi(n)})$。
该\omterm{def:b1:seq:subsequence}{子列}有界，故由波尔查诺–魏尔斯特拉斯定理
（见\cref{thm:b1:seq:bw}）可知它有一收敛的子\omterm{def:b1:seq:subsequence}{子列}，其极限
$\ell'$ 满足 $\abs{\ell' - \ell} \geq \varepsilon_0$
（对不等式取极限，见\cref{thm:b1:seq:order}）。但
$(u_n)$ 的子\omterm{def:b1:seq:subsequence}{子列}也是 $(u_n)$ 的收敛\omterm{def:b1:seq:subsequence}{子列}，
故由假设 $\ell' = \ell$，矛盾。
\end{solution}

\begin{solution}{exo:b1:seq:9}
设 $\eu = \frac pq$，$q \geq 1$。严格不等式 $a_q <
\eu < a_q + \frac{1}{q\,q!}$（因 $(a_n)$ 严格
递增、$(b_n)$ 严格递减），两边乘以 $q!$ 得
\[
q!\,a_q \;<\; q!\,\frac pq \;<\; q!\,a_q + \frac 1q \leq q!\,a_q + 1.
\]
又 $N = q!\,a_q = \sum_{k=0}^{q} \frac{q!}{k!}$ 为整数（每个
$\frac{q!}{k!}$ 都是整数之积，$k \leq q$），且
$q!\,\frac pq = (q-1)!\,p$ 亦为整数。由上式，整数
$(q-1)!\,p$ 严格位于 $N$ 与 $N + \frac 1q \leq N + 1$ 之间：
一个严格位于 $\intoo{N}{N+1}$ 内的整数，这是不可能的。因此
$\eu \notin \Q$.
\end{solution}

\begin{solution}{exo:b1:seq:10}
\begin{enumerate}
  \item 设 $\varepsilon > 0$ 及 $N$，使 $\abs{u_k - \ell} \leq
        \frac{\varepsilon}{2}$（当 $k > N$ 时）。对 $n > N$，有：
        \[
        \abs{c_n - \ell}
        = \Bigl|\frac{\sum_{k=1}^{n}(u_k - \ell)}{n}\Bigr|
        \leq \frac{\sum_{k=1}^{N} \abs{u_k - \ell}}{n}
        + \frac{n - N}{n}\cdot\frac{\varepsilon}{2}
        \leq \frac{C}{n} + \frac{\varepsilon}{2},
        \]
        其中 $C = \sum_{k=1}^N \abs{u_k - \ell}$ 为定值。当 $n$
        充分大时 $\frac Cn \leq \frac\varepsilon2$，于是 $\abs{c_n -
        \ell} \leq \varepsilon$.
  \item $u_n = (-1)^n$：发散，但 $c_n \to 0$（部分和
        有界（以 $1$ 为界），再除以 $n$）。
  \item 对数列 $v_n = u_{n+1} - u_n \to \ell$ 应用 (1)，其
        切萨罗均值为 $\frac{u_{n+1} - u_1}{n} \to \ell$
        （裂项相消），且 $\frac{u_{n+1}}{n} = \frac{u_{n+1} -
        u_1}{n} + \frac{u_1}{n} \to \ell$；重新调整下标
        （$\frac{u_n}{n} = \frac{u_n}{n-1}\cdot\frac{n-1}{n}$）可得
        $\frac{u_n}{n} \to \ell$.
\end{enumerate}
\end{solution}

\begin{solution}{exo:b1:seq:11}
设 $L = \inf_{n \geq 1} \frac{u_n}{n} \geq 0$，且 $\varepsilon >
0$。取 $m$ 使 $\frac{u_m}{m} \leq L + \varepsilon$。对每个 $n$
写成 $n = qm + r$，$0 \leq r < m$；由次可加性（迭代使用）得
$u_n \leq q\,u_m + u_r$，于是
\[
\frac{u_n}{n} \leq \frac{q m}{n}\cdot\frac{u_m}{m} + \frac{u_r}{n}
\leq \frac{u_m}{m} + \frac{\max(u_0, \dots, u_{m-1})}{n}
\leq L + \varepsilon + \frac{C_m}{n},
\]
利用 $qm \leq n$。当 $n$ 充分大时，$\frac{C_m}{n} \leq \varepsilon$：
于是 $L \leq \frac{u_n}{n} \leq L + 2\varepsilon$ 对所有充分大的 $n$ 成立，
即收敛于 $L$。
\end{solution}

\begin{solution}{exo:b1:seq:12}
$G = \Z + 2\pi\Z$ 是 $(\R, +)$ 的\omterm{def:b1:structures:subgroup}{子群}，它是稠密的：它不是
$\alpha\Z$，因为若 $1 = p\alpha$，$2\pi = q\alpha$，则 $2\pi =
\frac qp$ 为有理数——而 $\pi \notin \Q$（在本阶段不加证明地承认；其证明梗概见
\cref{ch:b1:integration}）。由 \cref{exo:b1:reals:9}，$G$
在 $\R$ 中稠密。

现设 $y \in \intcc{-1}{1}$，$\theta = \arcsin y$。由稠密性，
对每个 $\varepsilon > 0$，存在 $n \in \Z$，$k \in \Z$，使得
$\abs{(n + 2\pi k) - \theta} \leq \varepsilon$，即 $n$
在 $\varepsilon$ 范围内逼近 $\theta - 2\pi k$。由 $\sin$
的 $2\pi$ 周期性及 $1$-利普希茨性（$\abs{\sin a - \sin b} \leq
\abs{a - b}$，由中值不等式（见
\cref{ch:b1:derivative}）），
\[
\abs{\sin n - y} = \abs{\sin(n + 2\pi k) - \sin\theta}
\leq \abs{n + 2\pi k - \theta} \leq \varepsilon .
\]
注意：$n$ 取遍 $\Z$，但 $\sin(-n) = -\sin n$，而 $y$
在 $\intcc{-1}{1}$ 中任意，故只需取非负下标即可
（必要时将 $(n, y)$ 换为 $(-n, -y)$）。因此 $\{\sin n : n \in
\N\}$ 在 $\intcc{-1}{1}$ 中稠密。在\omterm{prop:b1:reals:intervals}{区间}上稠密的数列必有
\omterm{def:b1:seq:subsequence}{子列}趋于不同的值，故该数列发散。
\end{solution}

\begin{solution}{pb:b1:seq:1}
\textbf{1.} $\dfrac{n(n+1)/2}{n^2} = \dfrac{1 + 1/n}{2} \to
\dfrac12$，且 $\dfrac{n(n+1)(2n+1)/6}{n^3} = \dfrac{(1 +
1/n)(2 + 1/n)}{6} \to \dfrac13$.

\textbf{2.} \omterm{def:b1:reals:bounds}{上界}：$n$ 项中每一项都 $\leq \sqrt n$，故
$T_n \leq n\sqrt n$。下界：满足 $k > \frac n2$ 的项
个数至少为 $\frac n2$，而每一项都 $\geq \sqrt{n/2}$：
\[
T_n \geq \frac n2 \sqrt{\frac n2} = \frac{n^{3/2}}{2\sqrt 2} .
\]

\textbf{3.} 当 $k \geq N$ 时，由于 $b_{k+1} - b_k > 0$：
$m(b_{k+1} - b_k) \leq a_{k+1} - a_k \leq M(b_{k+1} - b_k)$.
对 $k = N, \dots, n - 1$ 求和，两边裂项相消：
\[
m\,(b_n - b_N) \leq a_n - a_N \leq M\,(b_n - b_N),
\]
两边除以 $b_n - b_N > 0$ 即得结论。

\textbf{4.} 由极限的定义，存在 $N$，使得 $\ell -
\varepsilon \leq \frac{a_{k+1} - a_k}{b_{k+1} - b_k} \leq \ell +
\varepsilon$ 对所有 $k \geq N$ 成立；由第 3 问，取 $m = \ell -
\varepsilon$，$M = \ell + \varepsilon$，该夹逼关系可
传递到 $\frac{a_n - a_N}{b_n - b_N}$ 上。

\textbf{5.} 展开恒等式的右边：
\[
\frac{a_N - \ell b_N}{b_n} + \frac{a_n - a_N}{b_n} -
\ell\,\frac{b_n - b_N}{b_n}
= \frac{a_n - \ell b_n}{b_n} = \frac{a_n}{b_n} - \ell .
\]
由第 4 问，乘积中的第二个因子
绝对值以 $\varepsilon$ 为界，且 $0 < 1 - \frac{b_N}{b_n}
\leq 1$（当 $n$ 充分大时），于是
\[
\Bigl|\frac{a_n}{b_n} - \ell\Bigr|
\leq \frac{\abs{a_N - \ell b_N}}{b_n} + \varepsilon
\leq 2\varepsilon
\]
只要 $b_n \geq \frac{\abs{a_N - \ell b_N}}{\varepsilon}$，
这一条件由于 $b_n \to +\infty$ 而最终成立。因此
$\frac{a_n}{b_n} \to \ell$，即切萨罗--施托尔茨定理。

\textbf{6.} 给定 $M$，取 $N$，使得 $\frac{a_{k+1} -
a_k}{b_{k+1} - b_k} \geq M$ 当 $k \geq N$ 时成立；由第 3 小题的
下半部分可得 $a_n - a_N \geq M(b_n - b_N)$，于是
\[
\frac{a_n}{b_n} \geq \frac{a_N}{b_n} +
M\Bigl(1 - \frac{b_N}{b_n}\Bigr) \longrightarrow M .
\]
从某项开始，$\frac{a_n}{b_n} \geq M - 1$；由于 $M$ 是
任意的，故 $\frac{a_n}{b_n} \to +\infty$。

\textbf{7.} 取 $b_n = n$，$a_n = u_1 + \dots + u_n$，则
增量商为 $u_{n+1} \to \ell$，于是切萨罗平均
$\frac{a_n}{n}$ 趋于 $\ell$，此即
\cref{exo:b1:seq:10} 的（1）。取 $a_n = u_n$，则增量商
为 $u_{n+1} - u_n$，即得（3）。反之：取 $a_n = (-1)^n$，
$b_n = n$，则 $\frac{a_n}{b_n} \to 0$，但 $a_{n+1} - a_n = \pm
2$ 交替变化，增量商没有极限。

\textbf{8.} 乘以共轭式：
\[
\bigl((1+h)^{3/2} - 1\bigr)\bigl((1+h)^{3/2} + 1\bigr)
= (1+h)^3 - 1 = 3h + 3h^2 + h^3 .
\]
取 $h = \frac1n$，则 $n\bigl((1 + \frac1n)^{3/2} - 1\bigr) =
\frac{3 + 3/n + 1/n^2}{(1 + 1/n)^{3/2} + 1}$，且 $1 \leq (1 +
\frac1n)^{3/2} \leq (1 + \frac1n)^2 \to 1$（夹逼），故
分母趋于 $2$，从而该式趋于 $\frac32$。现
用施托尔茨定理，取 $a_n = T_n$，$b_n = n^{3/2}$（严格递增，
$\to +\infty$):
\[
\frac{\sqrt{n+1}}{(n+1)^{3/2} - n^{3/2}}
= \frac{\sqrt{n+1}}{\sqrt n} \cdot
\frac{1}{n\bigl((1 + \frac1n)^{3/2} - 1\bigr)}
\longrightarrow 1 \cdot \frac{2}{3},
\]
从而 $T_n \sim \frac23\,n^{3/2}$。（第 2 小题的夹逼
只将常数夹在 $\intcc{0.35}{1}$ 中；施托尔茨定理则确定了它。）

\textbf{9.} 在 $u = \ln(1+t)$ 处用（G1）：$1 + t = \eu^{\ln(1+t)} \geq
1 + \ln(1+t)$，故 $\ln(1+t) \leq t$。取 $u =
-\frac{t}{1+t}$，由（G1）得 $\eu^{-t/(1+t)} \geq 1 - \frac{t}{1+t} =
\frac{1}{1+t} > 0$；两边取 $\ln$（单调递增），
$-\frac{t}{1+t} \geq -\ln(1+t)$，即 $\ln(1+t) \geq
\frac{t}{1+t}$。

\textbf{10.} $\ln$ 严格递增
（见\cref{prop:b1:functions:expln}），且 $\ln(2^k) = k\ln 2$ 无界，
故 $\ln n \to +\infty$。关于增量，取 $t =
\frac1n$，由第 9 小题，
\[
\frac 1{n+1} = \frac{1/n}{1 + 1/n} \leq
\ln\Bigl(1 + \frac1n\Bigr) \leq \frac1n
\quad\Longrightarrow\quad
\frac{n}{n+1} \leq \frac{1/(n+1)}{\ln(1 + 1/n)} \leq 1 ,
\]
故增量商 $\frac{H_{n+1} - H_n}{\ln(n+1) - \ln
n}$ 趋于 $1$，由施托尔茨定理得 $H_n \sim \ln n$。

\textbf{11.} 记 $x_n = \frac{u_{n+1}}{u_n} \to L$，$t_n =
\frac{x_n}{L} - 1 \to 0$。由第 9 小题，$\frac{t_n}{1 + t_n} \leq
\ln(1 + t_n) \leq t_n$，故 $\ln x_n - \ln L = \ln(1 + t_n) \to
0$（夹逼）。对 $(\ln x_k)$ 应用切萨罗平均（第 7 小题）：
\[
\frac1n \sum_{k=0}^{n-1} \ln x_k = \frac{\ln u_n - \ln u_0}{n}
\longrightarrow \ln L ,
\]
故 $\frac{\ln u_n}{n} \to \ln L$。记 $h_n = \frac{\ln u_n}{n}
- \ln L \to 0$，则 $u_n^{1/n} = L\,\eu^{h_n}$，由（G1）得
$1 + h_n \leq \eu^{h_n} \leq \frac{1}{1 - h_n}$（当 $h_n < 1$ 时），
故 $\eu^{h_n} \to 1$，从而 $u_n^{1/n} \to L$。应用：取 $u_n =
\binom{2n}{n}$，则
\[
\frac{u_{n+1}}{u_n} = \frac{(2n+1)(2n+2)}{(n+1)^2}
= \frac{2(2n+1)}{n+1} \longrightarrow 4 ,
\qquad\text{从而}\qquad
\binom{2n}{n}^{1/n} \to 4 .
\]

\textbf{12.} 当 $0 < x \leq 1$ 时，由（G2）得 $\sin x \geq x(1 -
\frac{x^2}{6}) \geq \frac{5x}{6} > 0$，且
\[
x - \sin x \geq \frac{x^3}{6} - \frac{x^5}{120}
= x^3\Bigl(\frac16 - \frac{x^2}{120}\Bigr)
\geq \frac{19}{120}\,x^3 > 0 :
\]
于是 $0 < \sin x < x$ 在 $\intoc{0}{1}$ 上成立。又 $u_1 = \sin u_0
\in \intcc{-1}{1}$。若 $u_1 = 0$，则 $u_n = 0$ 对所有 $n \geq 1$ 成立。
若 $u_1 \in \intoc{0}{1}$，由归纳法得 $0 < u_{n+1} = \sin u_n
< u_n \leq 1$，故 $(u_n)_{n\geq1}$ 严格递减，且
有下界 $0$，故收敛于某个 $\ell \in
\intco{0}{1}$（见\cref{thm:b1:seq:monotone}）。由和差化积公式
及（G2）可得 $\abs{\sin a - \sin b} = 2\abs{\cos
\frac{a+b}{2}}\,\abs{\sin\frac{a-b}{2}} \leq \abs{a - b}$，故
$u_{n+1} = \sin u_n \to \sin \ell$，于是 $\ell = \sin\ell$。若 $\ell
> 0$，则 $\sin\ell < \ell$，矛盾。故 $u_n \to 0$。

\textbf{13.} 将（G2）除以 $u_n > 0$：
\[
1 - \frac{u_n^2}{6} \leq \frac{\sin u_n}{u_n} \leq 1 -
\frac{u_n^2}{6} + \frac{u_n^4}{120} \leq 1 ,
\]
由 $u_n \to 0$ 夹逼得 $\frac{\sin u_n}{u_n} \to 1$。将
$x - \sin x$ 除以 $x^3$：
\[
\frac16 - \frac{u_n^2}{120} \leq \frac{u_n - \sin u_n}{u_n^3}
\leq \frac16 \longrightarrow \frac16 .
\]

\textbf{14.} 由于 $u_{n+1} = \sin u_n$，有
\[
w_n = \frac{u_n^2 - \sin^2 u_n}{u_n^2 \sin^2 u_n}
= \frac{(u_n - \sin u_n)(u_n + \sin u_n)}{u_n^2 \sin^2 u_n}
= \frac{u_n - \sin u_n}{u_n^3}\cdot
\frac{u_n + \sin u_n}{u_n}\cdot
\Bigl(\frac{u_n}{\sin u_n}\Bigr)^{2}
\]
（检查 $u_n$ 的幂次：$3 + 1 + (-4)$，与分母 $u_n^2$
及分子 $u_n^4$ 对照）。由第 13 小题，
三个因子分别趋于 $\frac16$、$2$、$1$，故 $w_n \to \frac13$。

\textbf{15.} 记 $v_n = \frac{1}{u_n^2}$，其增量为 $v_{n+1} -
v_n = w_n \to \frac13$，故 $\frac{v_n}{n} \to \frac13$，
由\cref{exo:b1:seq:10}~(3)得 $n u_n^2 \to 3$。于是
\[
\abs{\sqrt n\,u_n - \sqrt 3}
= \frac{\abs{n u_n^2 - 3}}{\sqrt n\,u_n + \sqrt 3}
\leq \frac{\abs{n u_n^2 - 3}}{\sqrt 3} \longrightarrow 0 :
\]
$\sqrt n\,u_n \to \sqrt 3$，即 $u_n \sim \sqrt{3/n}$。

\textbf{16.} 由 $n u_n^2 \to 3$，从某项起 $2 \leq n u_n^2
\leq 4$，即 $\sqrt{2/n} \leq u_n \leq 2/\sqrt n$。若 $u_n
\leq 10^{-2}$，且 $n$ 在该范围内，则 $2/n \leq 10^{-4}$：
$n \geq 20\,000$；而 $\sqrt{3/n} = 10^{-2}$ 对应 $n = 30\,000$。
\omterm{ex:b1:seq:heron}{海伦方法}每步使误差平方——有效数字的位数
加倍——因为在其不动点处，相关斜率的
绝对值 $< 1$（事实上此迭代是压缩的）。
此处 $\sin' 0 = \cos 0 = 1$：不动点是中性的，不存在
几何压缩，衰减速率由
首个非线性项 $-\frac{x^3}{6}$ 支配，因而是\omterm{def:b1:poly:def}{多项式}型。\omterm{ex:b1:seq:heron}{海伦方法}迭代一次
所得精度，超过正弦迭代一万次的
精度。

\textbf{17.} 对任意 $u_0$，$u_1 = \sin u_0 \in
\intcc{-1}{1}$。若 $u_1 = 0$，则数列
自第 $1$ 项起恒为零。若 $u_1 > 0$，第四部分的论证原样适用。若 $u_1 < 0$，令
$v_n = -u_n$。由 $\sin$ 的奇性，得 $v_{n+1} = -\sin u_n =
\sin(-u_n) = \sin v_n$，且 $v_1 \in \intoc{0}{1}$，故 $v_n \sim
\sqrt{3/n}$，即 $u_n \sim -\sqrt{3/n}$。总之
$\abs{u_n} \sim \sqrt{3/n}$（或数列最终为 $0$）：
下降是普遍的，唯有符号记得起点。

\textbf{18.} 先证 $\frac{u_{n+1}}{u_n} = 1 - \frac{u_n -
u_{n+1}}{u_n^2}\,u_n \to 1 - a \cdot 0 = 1$。于是
\[
\frac{1}{u_{n+1}} - \frac{1}{u_n}
= \frac{u_n - u_{n+1}}{u_n u_{n+1}}
= \frac{u_n - u_{n+1}}{u_n^2}\cdot\frac{u_n}{u_{n+1}}
\longrightarrow a \cdot 1 = a ,
\]
由\cref{exo:b1:seq:10}~(3)得 $\frac{1}{n u_n} \to a$，
即 $n u_n \to \frac1a$。

\textbf{19.} 令 $v_n = \frac{1}{u_n}$，则 $v_{n+1} = \frac{1 +
u_n}{u_n} = v_n + 1$，故 $v_n = v_0 + n$，且
\[
u_n = \frac{u_0}{1 + n u_0} , \qquad
n u_n = \frac{n u_0}{1 + n u_0} \longrightarrow 1 .
\]
验证引理：$u_n - u_{n+1} = \frac{u_n^2}{1 + u_n}$，故
$\frac{u_n - u_{n+1}}{u_n^2} = \frac{1}{1 + u_n} \to 1 = a$，而
第 18 小题预言 $n u_n \to 1$：两者完全吻合。

\textbf{20.} 由归纳法可知 $\eu^{-u} > 0$ 成立；
因 $\eu^{-u_n} < 1$（当 $u_n > 0$ 时），数列递减；故 $u_n \to
\ell \geq 0$（由\cref{thm:b1:seq:monotone}）。连续性桥梁：
记 $h_n = \ell - u_n \to 0$，则 $\eu^{-u_n} = \eu^{-\ell}
\eu^{h_n} \to \eu^{-\ell}$，由 (G1) 的夹逼 $1 + h_n \leq
\eu^{h_n} \leq \frac{1}{1 - h_n}$；故 $\ell = \ell\,
\eu^{-\ell}$。若 $\ell > 0$，则将推出 $\eu^{-\ell} = 1$，
矛盾，故 $\ell = 0$。对 $t > 0$，由 (G1) 得 $\eu^{-t} \geq 1 - t$
且 $\eu^{-t} \leq \frac{1}{1 + t}$，故
\[
\frac{1}{1 + t} \leq \frac{1 - \eu^{-t}}{t} \leq 1 .
\]
取 $t = u_n$，则 $\frac{u_n - u_{n+1}}{u_n^2} = \frac{1 -
\eu^{-u_n}}{u_n} \to 1$。由第 18 小题，取 $a = 1$，得 $n u_n \to
1$，故 $u_n \sim \frac1n$。

\textbf{21.} 与第 18 小题相同，$\frac{u_{n+1}}{u_n} = 1 -
\frac{u_n - u_{n+1}}{u_n^3}\,u_n^2 \to 1$。于是
\[
\frac{1}{u_{n+1}^2} - \frac{1}{u_n^2}
= \frac{(u_n - u_{n+1})(u_n + u_{n+1})}{u_n^2 u_{n+1}^2}
= \frac{u_n - u_{n+1}}{u_n^3}\cdot
\frac{u_n + u_{n+1}}{u_n}\cdot
\Bigl(\frac{u_n}{u_{n+1}}\Bigr)^2
\longrightarrow a \cdot 2 \cdot 1 = 2a ,
\]
由\cref{exo:b1:seq:10}~(3)得 $\frac{1}{n u_n^2} \to 2a$，
即 $n u_n^2 \to \frac{1}{2a}$。对正弦而言，$a = \frac16$
（第 13 小题）：$n u_n^2 \to 3$，恰为第四部分的结论。

\textbf{22.} 增量为：$a_{n+1} - a_n = 2^{-n} - 2^{-n-1} =
2^{-n-1} = b_{n+1} - b_n$，故增量之商
恒为 $1$。然而 $\frac{a_n}{b_n} = \frac{2 - 2^{-n}}{1 -
2^{-n}} \to 2 \neq 1$。第 5 小题的证明在
边界项处失效：$\frac{a_N - \ell b_N}{b_n} \to 0$ 需要 $b_n
\to +\infty$；此处（取 $\ell = 1$）$a_N - b_N = 1$，且 $b_n
\to 1$，故该项趋于 $1$——恰为剩余的
差距 $2 - 1$。

\textbf{23.} 应用施托尔茨定理，取 $A_n = \sum_{k=1}^n H_k$，$B_n = n\ln
n$：$B_{n+1} - B_n = \ln(n+1) + n\ln(1 + \frac1n) > 0$，且 $B_n
\to +\infty$。由第 9 小题，$\frac{n}{n+1} \leq n\ln(1 +
\frac1n) \leq 1$，故 $B_{n+1} - B_n = \ln(n+1) + \theta_n$，其中
$\frac12 \leq \theta_n \leq 1$。于是
\[
\frac{A_{n+1} - A_n}{B_{n+1} - B_n}
= \frac{H_{n+1}}{\ln(n+1)}\cdot
\frac{1}{1 + \theta_n/\ln(n+1)} \longrightarrow 1 \cdot 1 = 1
\]
（第 10 小题给出第一个因子的极限；$\theta_n$ 有界，
且 $\ln(n+1) \to \infty$ 给出第二个因子的极限）。施托尔茨定理由此得出：
$\sum_{k=1}^n H_k \sim n\ln n$。

\textbf{24.} 若 $u_n \to \ell > 0$：与第 11 小题相同，$\ln u_n
\to \ln\ell$（由第 9 小题对 $\ln\frac{u_n}{\ell}$ 的夹逼），故
切萨罗平均 $\frac1n\sum_{k=1}^n \ln u_k \to \ln\ell$，
由第 11 小题的指数桥梁得 $(u_1\cdots
u_n)^{1/n} = \exp\bigl(\frac1n\sum\ln u_k\bigr) \to \ell$。若
$u_n \to +\infty$：对任意 $M$，从某项起 $u_n \geq \eu^M$，故
$\ln u_n \geq M$，即 $\ln u_n \to +\infty$；由 $+\infty$ 情形的切萨罗
定理（第 6 小题，取 $b_n = n$）得 $\frac1n\sum \ln u_k \to
+\infty$，再由 (G1)（$\eu^s \geq 1 + s$）得几何平均
趋于 $+\infty$。取 $u_n = n$，则 $(n!)^{1/n} \to +\infty$。

\textbf{25.} (i) 完备性仅通过单调极限定理进入，
用以给出第 12 小题与第 20 小题中的极限；
切萨罗--施托尔茨定理本身是纯粹的
$\varepsilon$ 控制，在 $\Q$ 上成立。(ii) 施托尔茨定理将
$\lim \frac{a_n}{b_n}$ 替换为增量之商的 $\lim$，
恰如洛必达法则将 $\lim\frac fg$ 替换为
$\lim\frac{f'}{g'}$——其微分形式的对偶版本依赖于
在\cref{ch:b1:derivative}中的中值定理。(iii) 启发式地：若
$f(x) = x - a\,x^{p+1} + o(x^{p+1})$ 在中性不动点
$0$ 处，则 $\frac{1}{u_{n+1}^p} - \frac{1}{u_n^p} \to pa$，且
$u_n \sim (pan)^{-1/p}$：$p + 1$ 阶接触给出衰减
$n^{-1/p}$——图像相对于对角线越平坦，
下降越慢。(iv) 常数：(G2) 给出三次项
系数 $\frac16$；第 14 小题的因式分解使其加倍，
得望远镜式增量 $\frac13$；切萨罗将
$\frac{1}{u_n^2}$ 化为 $\frac n3$；取倒数并开方
即得 $\sqrt{3/n}$。
\end{solution}
